% Generated draft: requires mathematical and source review.
\documentclass[11pt]{article}
\usepackage[margin=25mm]{geometry}
\usepackage{fontspec}
\setmainfont{Latin Modern Roman}
\usepackage{amsmath,amssymb,mathtools,unicode-math}
\setmathfont{Latin Modern Math}
\usepackage{enumitem,needspace,xurl,hyperref}
\setlength{\parindent}{0pt}
\setlength{\parskip}{5pt}
\setlength{\emergencystretch}{3em}
\newcommand{\N}{\mathbb{N}}
\newcommand{\Z}{\mathbb{Z}}
\newcommand{\Q}{\mathbb{Q}}
\newcommand{\R}{\mathbb{R}}
\newcommand{\C}{\mathbb{C}}
\newcommand{\F}{\mathbb{F}}
\newcommand{\A}{\mathbb{A}}
\newcommand{\PP}{\mathbb{P}}
\newcommand{\E}{\mathbb{E}}
\newcommand{\eps}{\varepsilon}
\newcommand{\dd}{\,\mathrm{d}}
\begin{document}
\section*{UnsolvedMath ID: 30005877}
\textbf{Source problem code:} catalogue entry
\subsection{Definitions and mathematical statement}
Let \(G\) be a finite group with identity \(1_G\). For a generating subset \(A\subseteq G\), put \(S=A\cup A^{-1}\). Under the undirected Cayley graph convention, \(\Gamma(G,S)\) has vertex set \(G\), with an edge between \(g\) and \(gs\) for \(s\in S\). Its diameter is
\[
\operatorname{diam}\Gamma(G,S)
=\max_{g\in G}\min\{\ell\in\Z_{\geq 0}:g=s_1\cdots s_\ell,\ s_i\in S\},
\]
where the product of length zero is \(1_G\). Define the worst generating-set diameter by
\[
D(G)=\max_{\substack{A\subseteq G\\\langle A\rangle=G}}
\operatorname{diam}\Gamma(G,A\cup A^{-1}).
\]
The source does not specify whether its Cayley graphs are directed or whether generating sets are assumed symmetric; the convention above makes the undirected interpretation explicit.

For a positive integer \(n\), let \(\operatorname{Sym}_n\) be the group of permutations of \(\{1,\ldots,n\}\), and let \(\operatorname{Alt}_n\) be its subgroup of even permutations. For a field \(K\), write
\[
\operatorname{SL}_n(K)=\{M\in\operatorname{Mat}_{n\times n}(K):\det M=1\}.
\]
The source writes \(\operatorname{SL}_n\) without specifying a field, its size, or a quotient. The permutation and special linear groups mentioned as models are not all finite nonabelian simple groups; in particular, \(\operatorname{Alt}_n\) is nonabelian simple for \(n\geq5\), whereas no blanket simplicity assertion applies to the other named families.

\textbf{Source problem.} For finite nonabelian simple groups \(G\), develop a common framework for bounds on the diameter of every generating Cayley graph. Determine which growth arguments and proof notions for \(\operatorname{Sym}_n\) and \(\operatorname{Alt}_n\) can serve as models for \(\operatorname{SL}_n\) and, more broadly, linear algebraic groups, and which notions from the two settings can be formulated together. The supplied workshop question is a methodological research problem: it specifies neither a numerical bound nor a criterion that would uniquely determine a successful common framework. Its reference to ideas from ``the proof'' does not identify that proof.

\textbf{Qualified quantitative interpretation.} The supplied dated assessment describes the problem as a uniform Babai diameter and framework problem. A possible precise quantitative target under the undirected convention is
\[
\exists C>0\ \exists d>0\ \forall G\text{ finite nonabelian simple},\qquad
D(G)\leq C(\log |G|)^d,
\]
with constants independent of the group and the generating set. This formula is an interpretation of the assessment, not an explicit formula in the supplied workshop statement; it does not replace the separate request for transferable ideas and a common framework. The assessment also reports polylogarithmic bounds in bounded-rank Lie-type cases and unresolved alternating-group and unbounded-rank classical-group cases. These are supplied dated literature qualifications, not independently verified conclusions of this draft.
\subsection{Short English statement}
The problem asks for a common approach to bounding Cayley graph diameters of finite nonabelian simple groups, transferring ideas between permutation groups and linear algebraic groups. The source leaves the desired bound, graph convention, underlying field for \(\operatorname{SL}_n\), and meaning of a successful framework unspecified. A uniform polylogarithmic bound is recorded only as a qualified interpretation of the supplied literature assessment.
\subsection{Sources}
\begin{itemize}
\item Supplied local UnsolvedMath registry record, permanent ID \(30005877\), source problem code \texttt{catalogue entry}, title ``Uniform Diameter Bounds for Finite Simple Groups,'' in the Oberwolfach Reports Open Problems collection. The local registry is identified by the project as \texttt{data/unsolvedmath/problems.json}, with additional existing records in \texttt{catalog/problems.json}; no independent file inspection is claimed.
\item Supplied primary report citation: \emph{Mini-Workshop: Growth and Expansion in Groups} (2024), Oberwolfach Reports, DOI \texttt{10.4171/owr/2024/17}, pp.~1005--1040. \url{https://doi.org/10.4171/owr/2024/17}. The supplied extract associates the question with ``Strategies for growth in permutation groups'' and Harald Andr\'es Helfgott. The registry gives no separate proposer attribution or problem-specific page number.
\item Supplied literature source, labelled ``high-rank classical groups'': \url{https://link.springer.com/article/10.1007/s00222-021-01065-x}.
\item Supplied literature source, labelled ``unified growth perspective'': \url{https://arxiv.org/abs/1804.03049}.
\item Literature assessment date: \texttt{2026-08-22}. This is supplied dated curation derived from the OWR record, not independently verified literature review. The supplied status is ``open''; the assessment describes a still-open uniform Babai diameter and framework problem, reports polylogarithmic bounded-rank Lie-type cases, and identifies alternating groups and unbounded-rank classical groups as unresolved. The source expressly requests independent verification of time-sensitive status and cited claims.
\item The supplied statement status is \texttt{corrected\_verified}; its accompanying qualification says that only clear formula delimiters, notation, grammar, and typography were repaired, with mathematical content and scope preserved. This label does not establish mathematical correctness or independent source verification by this draft.
\end{itemize}
\end{document}
